\documentclass[11pt,letterpaper]{article}

\usepackage[top=1in, bottom=1in, left=1in, right=1in, headheight=14pt]{geometry}
\usepackage{fontspec}
\setmainfont{DejaVu Serif}
\setsansfont{DejaVu Sans}
\setmonofont{DejaVu Sans Mono}

\usepackage{xcolor}
\usepackage{titlesec}
\usepackage{fancyhdr}
\usepackage{enumitem}
\usepackage{hyperref}
\usepackage{booktabs}
\usepackage{tabularx}
\usepackage{longtable}
\usepackage{colortbl}
\usepackage{multirow}
\usepackage{array}
\usepackage{parskip}
\usepackage{tcolorbox}
\usepackage{graphicx}
\usepackage{lastpage}

% COLOR SCHEME — Industrial/Dark: Steel / Blood Burgundy / Aged Gold
\definecolor{primary}{HTML}{1C2A35}
\definecolor{secondary}{HTML}{7B1A1A}
\definecolor{tertiary}{HTML}{6B5614}
\definecolor{accent}{HTML}{A0522D}
\definecolor{lightprimary}{HTML}{E8ECF0}
\definecolor{lightsecondary}{HTML}{F5E6E6}
\definecolor{lighttertiary}{HTML}{F5F0E6}
\definecolor{rowgray}{HTML}{F5F5F5}
\definecolor{bordergray}{HTML}{BDBDBD}
\definecolor{subtlegray}{HTML}{757575}
\definecolor{darktext}{HTML}{212121}

% SECTION FORMATTING
\titleformat{\section}
  {\Large\bfseries\sffamily\color{primary}}
  {\thesection}{1em}{}
  [\vspace{-0.5em}\textcolor{bordergray}{\rule{\textwidth}{0.4pt}}]

\titleformat{\subsection}
  {\large\bfseries\sffamily\color{secondary}}
  {\thesubsection}{1em}{}

\titleformat{\subsubsection}
  {\normalsize\bfseries\sffamily\color{tertiary}}
  {\thesubsubsection}{1em}{}

\titlespacing*{\section}{0pt}{1.5em}{0.8em}
\titlespacing*{\subsection}{0pt}{1.2em}{0.5em}
\titlespacing*{\subsubsection}{0pt}{0.8em}{0.3em}

% CUSTOM ENVIRONMENTS
\newcommand{\stageheader}[2]{%
  \noindent\colorbox{#2}{\makebox[\textwidth][l]{%
    \textcolor{white}{\bfseries\sffamily\hspace{6pt}#1\hspace{6pt}}}}%
  \vspace{0.5em}\par%
}

\newtcolorbox{asciibox}{
  colback=rowgray, colframe=bordergray,
  arc=1pt, boxrule=0.5pt,
  left=6pt, right=6pt, top=4pt, bottom=4pt,
  fontupper=\small\ttfamily,
  breakable
}

\newtcolorbox{quotebox}{
  colback=lightprimary, colframe=primary,
  arc=2pt, boxrule=0.5pt,
  left=12pt, right=12pt, top=8pt, bottom=8pt,
  breakable
}

\newcommand{\metafield}[2]{\textbf{\sffamily #1} #2\par}

% TABLE HELPERS
\newcolumntype{L}[1]{>{\raggedright\arraybackslash}p{#1}}
\newcolumntype{C}[1]{>{\centering\arraybackslash}p{#1}}
\newcommand{\tableheader}[1]{\textbf{\sffamily\textcolor{white}{#1}}}

% HEADERS AND FOOTERS
\pagestyle{fancy}
\fancyhf{}
\fancyhead[L]{\small\sffamily\textcolor{subtlegray}{Post-Industrial Network}}
\fancyhead[R]{\small\sffamily\textcolor{subtlegray}{GSD Mission Package}}
\fancyfoot[C]{\small\sffamily\textcolor{subtlegray}{Page \thepage\ of \pageref{LastPage}}}
\renewcommand{\headrulewidth}{0.4pt}
\renewcommand{\footrulewidth}{0pt}

% HYPERLINKS
\hypersetup{
  colorlinks=true,
  linkcolor=secondary,
  urlcolor=secondary,
  pdfauthor={GSD Ecosystem},
  pdftitle={Post-Industrial Network -- Mission Package},
  pdfsubject={Vision to Mission Pipeline},
}

\begin{document}

% TITLE PAGE
\thispagestyle{empty}
\begin{center}
\vspace*{3cm}

{\small\sffamily\textcolor{primary}{\textbf{GSD RESEARCH MISSION PACKAGE}}}

\vspace{0.3cm}
\textcolor{bordergray}{\rule{8cm}{0.4pt}}

\vspace{1.5cm}
{\fontsize{28}{34}\selectfont\bfseries\sffamily\textcolor{primary}{The Post-Industrial\\[0.3em]Network}}

\vspace{0.8cm}
{\Large\sffamily\textcolor{secondary}{Throbbing Gristle, Nurse With Wound, Current 93,\\Death in June, Coil, and the Outlier Question}}

\vspace{0.5cm}
{\large\sffamily\itshape\textcolor{tertiary}{A Deep Research Study in Cultural Topology}}

\vspace{3cm}
{\sffamily\textcolor{accent}{Vision $\rightarrow$ Research $\rightarrow$ Mission}}\\[0.3em]
{\small\sffamily\textcolor{accent}{Full Pipeline Execution}}

\vspace{3cm}
{\small\sffamily\textcolor{subtlegray}{Date: March 26, 2026  |  Status: Ready for GSD Orchestrator}}\\[0.3em]
{\small\sffamily\textcolor{subtlegray}{Activation Profile: Fleet  |  Estimated: 14--18 context windows across 4--5 sessions}}
\end{center}
\newpage

% TABLE OF CONTENTS
\tableofcontents
\newpage

% ============================================================
% STAGE 1 — VISION DOCUMENT
% ============================================================
\stageheader{STAGE 1 --- VISION DOCUMENT}{primary}

\section{The Post-Industrial Network --- Vision Guide}

\metafield{Date:}{March 26, 2026}
\metafield{Status:}{Initial Vision / Conversation-Harvested}
\metafield{Depends on:}{gsd-coil-discography-study.md (Module 05 overshoot acknowledged)}
\metafield{Context:}{Companion study to the Coil deep research mission. Where the Coil study profiles a single
artist across 43 releases, this mission maps the ecosystem that made Coil possible — the relational
architecture of British post-industrial music from 1975 to the present.}

\subsection{Vision}

The Coil discography study, in its 479-line Module 05, revealed something the original scope could not
contain: Coil was never just a band. John Balance and Peter Christopherson were nodes in a network
— and you cannot understand the signal at a node without tracing the connections that feed it.
Jhonn Balance had been a member of Psychic TV, whose Genesis P-Orridge had invented industrial music in
the first place. Christopherson had been present at the origin of Throbbing Gristle, in the Hackney
studio they called the Death Factory. Current 93's David Tibet met Coil at The Equinox Event in 1983 and
collaborated for decades. Steven Stapleton of Nurse With Wound appeared on nearly every Current 93
release and became Coil's closest spiritual neighbor. Death in June's Douglas Pearce provided the folk
pivot that pulled the entire scene toward acoustic instrumentation. These are not footnotes to a Coil
biography. They are the context without which the biography makes no sense.

This study is organized around a core architectural insight from the Coil mission debrief: the value is
not in the nodes but in the spaces between them. What passes between artists — the shared personnel,
the collaborative recordings, the philosophical borrowings, the deliberate divergences — is where
post-industrial music's actual intelligence lives. The Nurse With Wound List, 291 entries printed on an
insert tucked into a debut album in 1979, is the most literal possible expression of this idea: it is
not a list of artists, it is a map of connections, a pre-internet index of the European experimental
underground. Its opening words are: ``categories strain, crack and sometimes break, under their burden —
step out of the space provided.''

The study must also grapple with the outlier question. Nine Inch Nails drew from the industrial palette
but routed it through arena-rock infrastructure, accumulating the cultural footprint of a major-label
phenomenon while the scene it borrowed from remained deliberately peripheral. Trent Reznor never called
his work industrial music; he borrowed techniques from Throbbing Gristle and Test Dept and then
delivered them to audiences who had never heard of either. Understanding how the industrial vocabulary
entered the mainstream — and what was lost and what was gained in that translation — completes the
network map by defining its outer boundary.

\subsection{Problem Statement}

\begin{enumerate}
  \item \textbf{The incomplete node profile problem.} Every existing account of this scene focuses on
    individual artists in isolation. No single document maps the full relational architecture — who
    collaborated with whom, when connections formed and broke, how philosophical positions influenced
    sonic choices across the network simultaneously.

  \item \textbf{The artifact-as-index problem.} The Nurse With Wound List is simultaneously a
    historical document, a collector's guide, and a structural map of 20 years of experimental
    music — but no research treatment handles all three functions simultaneously. Most treatments
    reduce it to a curiosity or a shopping list.

  \item \textbf{The divergence-without-comparison problem.} Current 93's move from industrial noise
    to apocalyptic folk, and Death in June's parallel transition from post-punk, are frequently
    described but rarely analyzed in comparison. The shared origin (Industrial Records milieu, Psychic
    TV adjacency, Throbbing Gristle as ur-text) makes the differences analytically meaningful.

  \item \textbf{The NIN framing problem.} Nine Inch Nails is either included in industrial surveys
    (distorting the center of gravity toward arena-rock aesthetics) or excluded (losing the only case
    study of what happens when the vocabulary crosses over). Neither choice is analytically satisfying.
    The study needs a structural position for NIN that acknowledges both its debt to the network and
    its fundamental difference from it.

  \item \textbf{The sensitivity and ideology problem.} Death in June in particular requires careful
    handling: Douglas Pearce's use of fascist imagery, collaboration with Boyd Rice, and the ideological
    drift of segments of the neofolk scene represent real ethical complexity. The study must address this
    directly rather than elide it.
\end{enumerate}

\subsection{Core Concept}

\textbf{Cultural topology over discography.} The unit of analysis is the \emph{network relationship},
not the individual artist. The study treats this scene as a directed graph: nodes are artists, edges
are collaborations, derivations, and philosophical influences. The Nurse With Wound List becomes
legible as an adjacency matrix. Current 93 and Death in June become case studies in divergent
evolution from a shared origin. NIN becomes a case study in what happens when a subgraph reaches
escape velocity into the mainstream.

\subsubsection{Network Map (Simplified)}

\begin{asciibox}
THE POST-INDUSTRIAL NETWORK — PRIMARY RELATIONSHIPS

  COUM Transmissions (1969-1976)
       |
       v
  THROBBING GRISTLE (1976-1981) -----> Industrial Records
  [Genesis P-Orridge, Peter Christopherson,         |
   Cosey Fanni Tutti, Chris Carter]                 |
       |                                            |
       +---------> PSYCHIC TV                       +---> Cabaret Voltaire
       |           [Genesis + Sleazy]                     SPK, Clock DVA
       |                |
       |                +--> COIL (1983-2004)
       |                     [Christopherson + Balance]
       |                           |
       |                           +<--> CURRENT 93 (1982-)
       |                           |    [David Tibet]
       |                           |         |
       |                    NWW List         +<--> NURSE WITH WOUND (1978-)
       |                    [291 entries]    |    [Steven Stapleton]
       |                                    |
       |                    DEATH IN JUNE --+
       |                    [Douglas P.]
       |                           |
       |                           v
       |                    NEOFOLK EMERGENCE
       |                    [Sol Invictus, Fire+Ice, etc.]
       |
       v
  NINE INCH NAILS (1988-) — OUTLIER NODE
  [Reznor: borrows TG techniques, routes through arena-rock]
  [The Downward Spiral: 4x platinum — escape velocity achieved]
\end{asciibox}

\subsubsection{Module Architecture}

\begin{asciibox}
MODULE STRUCTURE — 6 RESEARCH MODULES

  Module 01: The Origin Node
  |--> Throbbing Gristle / COUM Transmissions / Industrial Records
  |--> 1969-1981: From performance art to the genre-founding moment

  Module 02: The Cartographers
  |--> Nurse With Wound and the NWW List
  |--> The List as artifact, index, and structural map

  Module 03: The Divergent Path
  |--> Current 93 / David Tibet
  |--> Industrial noise --> apocalyptic folk --> Christian mysticism

  Module 04: The Folk Turn
  |--> Death in June and the neofolk emergence
  |--> Post-punk --> neofolk; ideology and ethics analysis

  Module 05: The Nexus (Cross-Reference Module)
  |--> Coil's position in the full network
  |--> Bridges to the companion Coil discography study

  Module 06: The Outlier Question
  |--> Nine Inch Nails and mainstream translation
  |--> What the vocabulary gains and loses in transit
\end{asciibox}

\subsection{Research Modules}

\subsubsection{Module 01 --- The Origin Node}

Throbbing Gristle are the ur-text. Industrial music can be traced to a precise moment of origin
(the ICA's \textit{Prostitution} exhibition, October 1976), a precise location (Hackney, London),
and a precise personnel (P-Orridge, Christopherson, Tutti, Carter). This module documents the
full arc from COUM Transmissions through Industrial Records' dissolution in 1981, including the
label's function as a distribution infrastructure for the emerging scene (Cabaret Voltaire, SPK,
Monte Cazazza, Clock DVA) and the significance of the band's self-termination as a methodological
statement.

\subsubsection{Module 02 --- The Cartographers}

The Nurse With Wound List is a 291-entry document printed on an insert tucked into a limited
edition debut album in January 1979. It maps 20 years of European experimental music — free
jazz, Krautrock, noise, outsider rock, musique concrète — without explanation, hierarchy, or
starting points. Steven Stapleton has since claimed some entries were invented; John Fothergill
denies this. The list opens with ``categories strain, crack and sometimes break.'' This module
analyzes the List as three simultaneous objects: historical artifact, collector's index, and
structural map of the scene's intellectual genealogy.

\subsubsection{Module 03 --- The Divergent Path}

David Tibet founded Current 93 in 1982, initially from within the Psychic TV orbit. By 1988,
he had pivoted to what he called ``apocalyptic folk'' — acoustic guitars, nursery rhymes,
Christian eschatology, Norse mythology. By the mid-1990s (*Thunder Perfect Mind*, widely
considered the apex of the project), Current 93 had traveled an extraordinary arc from
industrial noise to something closer to sacred music. Tibet's collaborations with Stapleton,
Balance, and Douglas Pearce make this arc the most richly documented divergence in the network.
This module traces the arc with album-by-album analysis of the key inflection points.

\subsubsection{Module 04 --- The Folk Turn}

Death in June emerged in 1981 from the post-punk band Crisis (explicitly left-wing and
anti-fascist). By 1984 Douglas Pearce had adopted acoustic guitars, European folk textures,
and — controversially — fascist imagery and collaboration with Boyd Rice. This module documents
the neofolk emergence with full analytical honesty: the music's power, the ideology's danger,
and the complex web of collaborations that tie Death in June to figures (Tibet, Stapleton,
Balance) with quite different political orientations. World Serpent Distribution is analyzed as
infrastructure that sustained the scene while also providing cover for its more troubling elements.

\subsubsection{Module 05 --- The Nexus}

Coil sits at the intersection of every other module. Christopherson was present at TG's origin.
Balance met Tibet at The Equinox Event. Stapleton and Coil were creative neighbors for two
decades. The first NIN/Coil collaboration (the \textit{Fixed} EP, 1992) is the single moment
where the network's interior and the mainstream outlier made direct contact. This module is
explicitly a bridge document — it cross-references the companion Coil discography study and
provides the analytical frame for how Coil synthesized inputs from every other node in the graph.

\subsubsection{Module 06 --- The Outlier Question}

Reznor borrowed Throbbing Gristle and Test Dept techniques and routed them through synth-pop
hooks, major label infrastructure, and arena-rock presentation. \textit{Pretty Hate Machine}
(1989) went triple platinum. \textit{The Downward Spiral} (1994) debuted at \#2 on the
Billboard 200 and went quadruple platinum. Reznor has explicitly stated he does not consider
his work ``industrial music.'' This module analyzes NIN as an outlier case: what the industrial
vocabulary does when it reaches escape velocity, why the originating network is structurally
incapable of replicating that trajectory, and what the Coil/NIN collaboration on \textit{Fixed}
(1992) reveals about the relationship between the two worlds.

\subsection{Chipset Configuration}

\begin{asciibox}
# chipset.yaml — Post-Industrial Network Study
mission_id: post-industrial-network-v1
activation_profile: fleet
model_allocation:
  opus: 30%   # Synthesis, ideology analysis, network topology judgment
  sonnet: 60% # Module surveys, discography tables, document assembly
  haiku: 10%  # Shared schemas, templates, scaffolding

craft_agents:
  - CRAFT-HIST: Music history and cultural context specialist
  - CRAFT-TOPO: Network topology and relationship mapping specialist
  - CRAFT-ETHI: Ideological/ethical sensitivity specialist (Module 04 primary)
  - CRAFT-ARCH: Cross-module architecture and synthesis specialist

parallel_capacity: 3
wave_count: 4
estimated_tokens: 280000-340000
\end{asciibox}

\subsection{Scope Boundaries}

\subsubsection{In Scope (v1.0)}

\begin{itemize}[noitemsep]
  \item Full historical documentation of Throbbing Gristle / COUM Transmissions arc (1969--1981)
  \item The Nurse With Wound List as primary artifact: composition, influence, legacy
  \item Current 93 discographic arc from \textit{LAShTAL} (1983) through \textit{Thunder Perfect Mind} (1993) with key later works
  \item Death in June from formation (1981) through neofolk codification (\textit{But, What Ends When the Symbols Shatter?}, 1992)
  \item Coil's network position as synthesis node (with cross-reference to companion study)
  \item Nine Inch Nails as outlier: \textit{Pretty Hate Machine}, \textit{Broken EP}, \textit{The Downward Spiral}, and the \textit{Fixed}/Coil collaboration
  \item Network relationship mapping: collaborative credits, shared labels, ideological borrowings
  \item Ethical analysis of fascist imagery and neofolk's ideological drift
  \item \textit{England's Hidden Reverse} (David Keenan, 2003) as primary secondary source
  \item World Serpent Distribution as network infrastructure
\end{itemize}

\subsubsection{Out of Scope}

\begin{itemize}[noitemsep]
  \item Full discographic treatment of Nurse With Wound (40+ releases beyond the List itself)
  \item Boyd Rice / NON — referenced but not profiled
  \item Cabaret Voltaire, SPK, Clock DVA — referenced as Industrial Records context but not profiled
  \item Sol Invictus, Fire and Ice, Forseti — referenced as neofolk second-generation but not profiled
  \item Genesis P-Orridge biography beyond TG/Psychic TV arc
  \item Electronic Body Music (EBM) and its divergence from industrial (separate study warranted)
  \item Post-2000 NIN output beyond network-relevant cross-references
\end{itemize}

\subsection{Success Criteria}

\begin{enumerate}
  \item \textbf{Origin documentation:} Throbbing Gristle arc from COUM Transmissions through
    Industrial Records dissolution is fully documented with key dates, personnel, and
    contextual significance.
  \item \textbf{NWW List analysis:} The Nurse With Wound List is analyzed at all three levels
    (artifact, index, structural map) with geographic and genre distribution of its 291 entries.
  \item \textbf{Current 93 arc:} At least eight key albums are individually analyzed marking
    the industrial-to-folk transition, with David Tibet's theological evolution traced.
  \item \textbf{Death in June arc:} The post-punk-to-neofolk transition is documented with
    full analysis of the ideological complexity and its implications for the scene.
  \item \textbf{Network topology:} A complete relationship map documents at least 20 distinct
    artist-to-artist edges with evidence (collaborative credits, shared personnel, documented
    influence).
  \item \textbf{Nexus analysis:} Coil's structural position in the network is analytically
    established with citations to the companion discography study.
  \item \textbf{NIN outlier framing:} Nine Inch Nails is positioned as an outlier node with
    clear analysis of what distinguished its trajectory from the originating network.
  \item \textbf{Coil/NIN contact point:} The \textit{Fixed} EP collaboration (1992) and
    subsequent \textit{Downward Spiral} remixes are analyzed as the network's primary
    mainstream contact moment.
  \item \textbf{Ethical analysis:} Death in June's ideological complexity is addressed with
    scholarly honesty — neither elided nor used to invalidate the music.
  \item \textbf{Sensitivity compliance:} All claims about political ideology are attributed
    to specific documented sources; no generalizations about entire genre communities.
  \item \textbf{Self-containment:} A reader with no prior knowledge of the scene can follow
    the full document without external reference.
  \item \textbf{Companion integration:} Module 05 successfully bridges to the companion
    Coil discography study with explicit cross-references to its module structure.
\end{enumerate}

\subsection{Relationship to Other Documents}

\begin{center}
\rowcolors{2}{rowgray}{white}
\begin{tabularx}{\textwidth}{L{4cm}L{3cm}X}
\toprule
\rowcolor{primary}
\tableheader{Document} & \tableheader{Relationship} & \tableheader{Notes} \\
\midrule
Coil Discography Study & Peer / Parent & This mission is the companion network study that the Coil study's Module 05 overshoot revealed was needed \\
gsd-amiga-vision & Philosophical & Amiga Principle applies: specialized nodes + architectural leverage over brute-force coverage \\
gsd-bbs-educational-pack & Structural & BBS pack's information-network model parallels the post-industrial scene's pre-internet distribution infrastructure \\
The Space Between (textbook) & Through-line & Network topology as philosophical principle; meaning in connections over nodes \\
\bottomrule
\end{tabularx}
\end{center}

\subsection{The Through-Line}

This is a study about what happens in the spaces between artists. The Nurse With Wound List
exists not to define a canon but to map a field of possibility — and its opening declaration
(``categories strain, crack and sometimes break'') is both a methodological instruction and a
philosophical position. The artists in this network were not building a genre. They were
maintaining pressure against the stabilization of genre — keeping the vocabulary unstable,
contested, open.

The GSD ecosystem's own philosophical foundation — that architectural leverage comes from
specialized nodes connected with integrity, not from brute-force coverage — maps onto this scene
with unusual precision. Throbbing Gristle invented industrial music and then deliberately
terminated themselves so it could not be reduced to a style. Coil never released a consistent
album. Current 93 changed its sound every few years. The network sustained itself through
constant self-disruption. That is the Amiga Principle operating at a cultural scale: small,
principled components producing staggering complexity through faithful iteration, with no node
large enough to calcify the whole.

Nine Inch Nails matters to this study not because it belongs to the network, but because it
defines the network's outer boundary — the point at which the vocabulary detaches from the
infrastructure that generated it. Understanding the boundary is how you understand the shape.

% ============================================================
% STAGE 2 — RESEARCH REFERENCE
% ============================================================
\newpage
\stageheader{STAGE 2 --- RESEARCH REFERENCE}{secondary}

\section{The Post-Industrial Network --- Research Reference}

\subsection{How to Use This Document}

This research reference provides module-by-module findings synthesized from primary source research.
All claims are sourced. Agents executing this mission should use this reference as the primary
evidence base, supplementing with targeted web research for gap-filling only.

\textbf{Key source organizations and texts:}
\begin{itemize}[noitemsep]
  \item \textit{England's Hidden Reverse} — David Keenan (2003; revised 2016). The definitive history
    of Coil, Current 93, and Nurse With Wound.
  \item Industrial Records / Mute Records — primary release documentation
  \item Wikipedia (music articles with citations) — biography and discography verification
  \item Rate Your Music, Discogs — release documentation and collaborative credits
  \item Stereogum (NWW List at 40 feature, 2019) — List analysis and anniversary coverage
  \item The Wire magazine — primary critical coverage of the scene since the 1980s
\end{itemize}

\subsection{Module 01 Reference --- The Origin Node}

\subsubsection{COUM Transmissions (1969--1976)}

COUM Transmissions was a Fluxus-influenced performance art collective formed in Hull in 1969 by
Genesis P-Orridge (born Neil Megson). The group relocated to Hackney, London in 1973, where
P-Orridge and Cosey Fanni Tutti met Chris Carter (then a sound recordist for television) and
Peter Christopherson (then a member of the graphic design collective Hipgnosis). The four built
a recording studio in Hackney, which they named the Death Factory.

The final COUM Transmissions performance — the \textit{Prostitution} exhibition at London's
Institute of Contemporary Arts, October 1976 — was simultaneously the public debut of
Throbbing Gristle. The exhibition generated parliamentary condemnation: Conservative MP Nicholas
Fairbairn declared the group ``wreckers of civilisation.'' The scandal confirmed the group's
central method: exposure as a form of revelation.

\subsubsection{Throbbing Gristle (1976--1981)}

\begin{center}
\rowcolors{2}{rowgray}{white}
\begin{tabularx}{\textwidth}{L{3.5cm}L{2.5cm}X}
\toprule
\rowcolor{primary}
\tableheader{Album / Release} & \tableheader{Year} & \tableheader{Significance} \\
\midrule
The Second Annual Report & 1977 & Debut; framed as documentation, not art. Ltd. 786 copies on Industrial Records \\
D.o.A: The Third and Final Report & 1978 & Expanded sonic palette; synthesizer-led tracks \\
20 Jazz Funk Greats & 1979 & Ironic title belied by the content; experimental peak \\
Heathen Earth & 1980 & Fully live studio performance; unique document of the band's process \\
Mission Terminated postcard & 1981 & Self-dissolution announced via mailed postcard. ``The mission is terminated.'' \\
\bottomrule
\end{tabularx}
\end{center}

Industrial Records (founded 1976) functioned as both label and conceptual artwork. Its slogan —
``Industrial Music for Industrial People'' — was coined with Monte Cazazza to reference both
the sonic palette and the mass unemployment reshaping British manufacturing. The label distributed
work by Cabaret Voltaire, SPK, Clock DVA, William S. Burroughs, and Monte Cazazza. Every TG
concert was recorded and made available via mail order — an infrastructure that prefigured
the internet's long-tail distribution model by two decades.

\textbf{Post-TG dispersal:} P-Orridge and Christopherson formed Psychic TV. Tutti and Carter
formed Chris and Cosey. These divergent paths established the two primary downstream branches
of the network: the noise-ritual-occult branch (Psychic TV $\to$ Coil) and the electronic-body
branch (Chris and Cosey). John Balance, a Psychic TV member, later partnered with Christopherson
to form Coil — closing a loop between the network's first and second generation.

\subsection{Module 02 Reference --- The Cartographers}

\subsubsection{Nurse With Wound: Formation and Method}

Nurse With Wound formed in 1978 when visual artist Steven Stapleton, working as a sign painter
at a London recording studio, improvised a fictional band to secure cheap studio time. He recruited
John Fothergill and Heman Pathak, gathered ``cheap guitars, synthesizers and other found objects,''
and recorded the debut in six hours. The album was issued in January 1979 on the group's own
United Dairies label (named after Rolf-Ulrich Kaiser's Krautrock label Ohr) in a limited pressing
of 500 copies.

\subsubsection{The Nurse With Wound List}

\begin{center}
\rowcolors{2}{rowgray}{white}
\begin{tabularx}{\textwidth}{L{3cm}C{2cm}X}
\toprule
\rowcolor{primary}
\tableheader{Parameter} & \tableheader{Value} & \tableheader{Notes} \\
\midrule
Initial entry count & 236 & Included with \textit{Chance Meeting} (1979) \\
Expanded entry count & 291 & Revised for \textit{To the Quiet Men from a Tiny Girl} (1980) \\
Geographic range & International & Strong European core; French, German, Italian, UK, Japanese, American entries \\
Genre spread & Broad & Free jazz, Krautrock, noise, musique concrète, progressive rock, psychedelic, outsider \\
Named canonical entries & Multiple & Velvet Underground, Captain Beefheart, Yoko Ono, Can, La Monte Young, Frank Zappa \\
Finders Keepers compilation & 2019 & \textit{Strain, Crack \& Break} Vol. 1 (French entries); Vol. 2 (German entries) \\
\bottomrule
\end{tabularx}
\end{center}

The List's opening declaration — ``categories strain, crack and sometimes break, under their
burden — step out of the space provided'' — establishes its methodological position. It is not
a canon. It is a refusal of canon. The entries range from the commercially available (Velvet
Underground) to the genuinely impossible-to-find (Sphinx Tush: a single song in two live versions,
1970). Stapleton has claimed some entries were invented; Fothergill denies this. The ambiguity is
structural: the List was designed to be followed up, not simply consumed.

Stereogum's 40th anniversary assessment (2019) noted that the NWW List prefigured the internet's
discovery infrastructure by two decades — a pre-algorithmic recommendation engine encoded in
a vinyl insert. Finders Keepers Records responded to the anniversary by launching the
\textit{Strain, Crack \& Break} compilation series, drawing exclusively from List entries.

\subsubsection{Nurse With Wound's Evolution}

By 1981, Stapleton was the sole remaining NWW member. He considers \textit{Homotopy to Marie}
(1982) the first proper NWW release. The project subsequently evolved from industrial noise
toward surrealism, dada, dark ambient, and drone — drawing on Stapleton's visual art background
and his fondness for Dadaism and absurdist humor. Over 40 full-length titles followed. Stapleton
appeared on nearly every Current 93 release until 2010, making him the most consistent node
in the network's inner circle.

\subsection{Module 03 Reference --- The Divergent Path}

\subsubsection{Current 93: Arc Overview}

\begin{center}
\rowcolors{2}{rowgray}{white}
\begin{tabularx}{\textwidth}{L{3.5cm}L{2.5cm}X}
\toprule
\rowcolor{primary}
\tableheader{Album} & \tableheader{Year} & \tableheader{Era / Significance} \\
\midrule
LAShTAL & 1983 & Industrial noise debut; Tibet describes as ``invocation of Malkunofath'' \\
Nature Unveiled & 1984 & Dark ambient; co-released with NWW collaboration \\
Dogs Blood Rising & 1984 & Harsh electronics; early Coil collaboration (Balance co-founders) \\
Swastikas for Noddy & 1988 & First full apocalyptic folk album; acoustic guitars enter \\
Earth Covers Earth & 1988 & Transition confirmed; folk textures dominant \\
Thunder Perfect Mind & 1992 & Critical peak; accessible entry point; last album with Douglas Pearce \\
Of Ruine, Or Some Blazing Starre & 1994 & Post-Pearce; deeper into Christian mysticism \\
I Am the Last of All the Field & 2001 & Mature neofolk; Tibet now identifies explicitly as Christian \\
\bottomrule
\end{tabularx}
\end{center}

\subsubsection{The Equinox Event (1983)}

The Equinox Event was a festival of post-industrial acts at which Tibet met Steven Stapleton.
Current 93 performed under the name Dogs Blood Order. John Balance (Coil's co-founder) was
in the original Current 93 lineup alongside Fritz Catlin of 23 Skidoo. Tibet's invitation to
Balance to leave Psychic TV and join Current 93 full-time was refused — a refusal that
indirectly led to Coil's formation and shaped the network's topology for the next two decades.

\subsubsection{Tibet's Theological Evolution}

Tibet took his name from Genesis P-Orridge, reflecting early interest in Tibetan Buddhism and
Thelema. Current 93's name derives from Aleister Crowley's ``93 Current'' in Thelema. The
trajectory moved from Crowleyan occultism through pagan Norse mythology (period of living with
Death in June's Douglas P. and Freya Aswynn) to what Tibet now describes as Christian faith,
specifically Christian mysticism and eschatology. Tibet completed a Master's degree in Coptic
Language and Grammar in 2009. His archives through 2013 are held at Cornell University.

\subsection{Module 04 Reference --- The Folk Turn}

\subsubsection{Death in June: Formation and Arc}

Death in June formed in 1981. Douglas Pearce and Tony Wakeford had been members of the
left-wing anti-fascist punk band Crisis (formed 1977). The initial Death in June sound retained
post-punk structures with martial drums and abrasive electronics. The folk turn began with
\textit{Burial} (1984): acoustic guitars, references to European history, and the introduction
of the Death's Head imagery that would define the group's visual identity and generate decades
of controversy.

\begin{center}
\rowcolors{2}{rowgray}{white}
\begin{tabularx}{\textwidth}{L{4cm}L{2cm}X}
\toprule
\rowcolor{primary}
\tableheader{Album} & \tableheader{Year} & \tableheader{Significance} \\
\midrule
The Guilty Have No Pride & 1983 & Post-punk debut \\
Burial & 1984 & Folk turn begins; acoustic guitars; first fascist imagery \\
Brown Book & 1987 & Consolidation of neofolk sound; introduces Ian Read \\
But, What Ends When the Symbols Shatter? & 1992 & Neofolk codification; widely cited as definitive statement \\
\bottomrule
\end{tabularx}
\end{center}

\subsubsection{Ideological Complexity}

Death in June's adoption of fascist imagery — SS Totenkopf skulls, runes, references to
Ernst Röhm and the Night of the Long Knives — has generated sustained controversy. The
analytical complexity is genuine: Pearce is openly gay and has collaborated extensively
with Jewish musicians; the Death in June website previously displayed an Israeli flag; the
band performed in Israel in 2004. Collaborator Boyd Rice (NON), however, has documented
connections to the Church of Satan and neo-fascist movements. David Tibet, a longtime
collaborator, has publicly distanced himself from Pearce over these associations.

World Serpent Distribution (founded 1991 by Pearce) distributed neofolk and post-industrial
releases across the UK through the late 1990s, functioning as the scene's primary commercial
infrastructure — and inadvertently creating a shared distribution channel for both its
more ethically complex and more ethically straightforward participants.

\subsubsection{Neofolk Emergence}

\begin{center}
\rowcolors{2}{rowgray}{white}
\begin{tabularx}{\textwidth}{L{3cm}L{3cm}X}
\toprule
\rowcolor{primary}
\tableheader{Artist} & \tableheader{Period} & \tableheader{Role in Network} \\
\midrule
Death in June & 1981-- & Originator; acoustic folk pivot; controversial imagery \\
Current 93 & 1982-- & Parallel originator; ``apocalyptic folk'' term coined by Tibet \\
Sol Invictus & 1987-- & Tony Wakeford (ex-DiJ); less ideologically fraught variant \\
Fire and Ice & 1987-- & Ian Read; introduced via \textit{Brown Book} \\
Swans & 1982-- & Parallel trajectory; \textit{Love of Life} (1992) embraces neofolk \\
\bottomrule
\end{tabularx}
\end{center}

\subsection{Module 05 Reference --- The Nexus}

Coil's structural position in the network is analyzed in the companion discography study.
Key cross-reference points for this mission:

\begin{itemize}[noitemsep]
  \item Christopherson's direct presence at TG's origin (TG founding member)
  \item Balance's early Current 93 membership (Dogs Blood Order period, 1982--1984)
  \item Stapleton's near-permanent presence on Current 93 releases = NWW/Coil crossover
  \item The Equinox Event (1983) as the network's founding convergence moment
  \item \textit{The Snow EP} (1993) — Coil's closest gesture toward accessibility
  \item \textit{Fixed} (1992) — First Coil/NIN collaboration; the network's primary mainstream contact
  \item \textit{Closer (Precursor)} on \textit{Closer to God} (1994) — Coil's finest mainstream-adjacent work
\end{itemize}

\subsection{Module 06 Reference --- The Outlier Question}

\subsubsection{NIN: Industrial Vocabulary, Arena Infrastructure}

Nine Inch Nails was formed by Trent Reznor in 1988 in Cleveland, Ohio. Reznor has explicitly
stated he does not consider his work ``industrial music'' but acknowledges borrowing techniques
from Throbbing Gristle and Test Dept. \textit{Pretty Hate Machine} (1989) deployed industrial
vocabulary (drum machines, synthesizer distortion, noise textures) alongside synth-pop hooks
and melodic vocals — a combination that had no precedent in the originating network.

\begin{center}
\rowcolors{2}{rowgray}{white}
\begin{tabularx}{\textwidth}{L{4cm}L{2cm}X}
\toprule
\rowcolor{primary}
\tableheader{Release} & \tableheader{Year} & \tableheader{Network Relevance} \\
\midrule
Pretty Hate Machine & 1989 & Vocabulary borrowed; independent release; 3$\times$ platinum \\
Fixed (EP) & 1992 & First Coil collaboration; J.G. Thirlwell also contributes \\
The Downward Spiral & 1994 & Recorded at Tate House (``Le Pig''); \#2 Billboard 200; 4$\times$ platinum \\
Closer to God (EP) & 1994 & Coil's ``Closer (Precursor)'' — the network's most-heard artifact \\
Further Down the Spiral & 1995 & Coil, Aphex Twin, Rick Rubin contribute; network-mainstream dialogue peaks \\
\bottomrule
\end{tabularx}
\end{center}

\subsubsection{The Structural Distinction}

The originating network sustained itself through deliberate peripheral positioning: limited
pressings, anti-commercial aesthetics, self-dissolution (TG's postcard), refusal of genre
labeling. NIN's trajectory was structurally opposite: major label distribution, MTV placement,
Lollapalooza touring, Grammy nominations. Coil's \textit{Closer (Precursor)} reached more ears
via the \textit{Se7en} opening credits (where it appeared as David Fincher's edit of
\textit{Closer to God}) than via any Coil release — an irony that encodes the outlier
relationship precisely.

The Reznor question does not resolve into ``NIN belongs to the network'' or ``NIN doesn't.''
It resolves into: the industrial vocabulary can achieve escape velocity, but the infrastructure
that generates the vocabulary is structurally incapable of following it into the mainstream.
The network's anti-commercial architecture is not a failure mode; it is the mechanism by which
the vocabulary stays alive and potent at the periphery.

\subsection{Safety and Sensitivity Considerations}

\begin{center}
\rowcolors{2}{rowgray}{white}
\begin{tabularx}{\textwidth}{L{3cm}L{2.5cm}X}
\toprule
\rowcolor{primary}
\tableheader{Sensitivity Area} & \tableheader{Boundary Type} & \tableheader{Required Handling} \\
\midrule
Fascist imagery (Death in June) & GATE & Document with source attribution; no interpretive generalization to full genre \\
Political ideology (neofolk right-wing drift) & GATE & Distinguish individual artists; cite specific documented associations \\
Occult/esoteric content & ANNOTATE & Present as cultural context; no endorsement or dismissal \\
Genesis P-Orridge biography & ANNOTATE & Use correct pronouns per P-Orridge's own statements (she/her in later life) \\
Controversial collaboration credits & GATE & Attribute all political claims to named sources only \\
\bottomrule
\end{tabularx}
\end{center}

\subsection{Source Bibliography}

\subsubsection{Primary Reference Texts}
\begin{itemize}[noitemsep]
  \item Keenan, David. \textit{England's Hidden Reverse: A Secret History of the Esoteric Underground.} Strange Attractor Press, 2003; revised 2016.
  \item Ford, Simon. \textit{Wreckers of Civilisation: The Story of COUM Transmissions and Throbbing Gristle.} Black Dog Publishing, 1999.
  \item Tutti, Cosey Fanni. \textit{Art Sex Music.} Faber and Faber, 2017.
\end{itemize}

\subsubsection{Discographic and Archival Sources}
\begin{itemize}[noitemsep]
  \item Industrial Records catalog documentation (Mute Records reissue liner notes)
  \item Discogs artist pages: Throbbing Gristle, Nurse With Wound, Current 93, Death in June, Coil, Nine Inch Nails
  \item Rate Your Music artist profiles and discographies (all six primary artists)
  \item NWW List Wikipedia article (February 2026 revision)
  \item Nurse With Wound official discography: ultimathulerecords.com
\end{itemize}

\subsubsection{Journalism and Critical Accounts}
\begin{itemize}[noitemsep]
  \item Stereogum: ``The Nurse With Wound List at 40: A Beginner's Guide.'' 2019.
  \item The Wire magazine: Ongoing coverage, 1982--present.
  \item Nine Inch Nails official discography and release notes: nin.com
  \item Wikipedia: Current 93, Throbbing Gristle, Death in June, Neofolk (all accessed March 2026)
  \item Cornell University Library: David Tibet Archives (through 2013)
\end{itemize}

% ============================================================
% STAGE 3 — MISSION PACKAGE
% ============================================================
\newpage
\stageheader{STAGE 3 --- MISSION PACKAGE}{tertiary}

\section{Milestone Specification}

\subsection{Mission Objective}

Produce a comprehensive six-module research document mapping the post-industrial music network
from its origin in COUM Transmissions (1969) through the neofolk emergence and the NIN outlier
case, providing analytical tools (network topology, ideological analysis, outlier framing) that
complement the companion Coil discography study and stand independently as a reference for the
British experimental underground.

\subsection{Architecture Overview}

\begin{asciibox}
WAVE ARCHITECTURE — Post-Industrial Network Study

Wave 0: FOUNDATION
  [HAI] Schema scaffolding, source index, shared templates
  Sequential | 1 session | ~15,000 tokens

Wave 1: PARALLEL SURVEYS
  Track A [SON]  Track B [OPU]   Track C [SON]
  Module 01      Module 02       Modules 03+04
  TG/COUM/IR     NWW+List        C93 + DiJ
  Parallel | 2 sessions | ~90,000 tokens per track

Wave 2: SYNTHESIS
  [OPU] Network topology construction | Outlier analysis
  [OPU] Sensitivity review (Module 04 ideological content)
  Sequential | 1-2 sessions | ~60,000 tokens

Wave 3: PUBLICATION + VERIFICATION
  [SON] Document assembly + formatting
  [SON] Source verification pass
  [SON] Module 05 cross-reference integration
  Sequential | 1 session | ~40,000 tokens
\end{asciibox}

\subsection{Deliverables}

\begin{center}
\rowcolors{2}{rowgray}{white}
\begin{tabularx}{\textwidth}{C{0.5cm}L{4cm}L{4cm}X}
\toprule
\rowcolor{primary}
\tableheader{\#} & \tableheader{Deliverable} & \tableheader{Acceptance Criteria} & \tableheader{Spec Location} \\
\midrule
1 & Module 01: Origin Node & TG arc fully documented; Industrial Records analyzed & Stage 2, \S2.2 \\
2 & Module 02: Cartographers & NWW List analyzed at all three levels; entry distribution mapped & Stage 2, \S2.3 \\
3 & Module 03: Divergent Path & Current 93 arc traced with 8+ album analyses & Stage 2, \S2.4 \\
4 & Module 04: Folk Turn & DiJ arc + neofolk emergence + ideological analysis & Stage 2, \S2.5 \\
5 & Module 05: Nexus & Coil's position mapped; companion study cross-referenced & Stage 2, \S2.6 \\
6 & Module 06: Outlier & NIN outlier analysis; Coil/NIN contact documented & Stage 2, \S2.7 \\
7 & Network topology diagram & 20+ edges documented with evidence & Stage 1, \S1.3 \\
8 & Full bibliography & All claims sourced; zero entertainment media & Stage 2, \S2.9 \\
\bottomrule
\end{tabularx}
\end{center}

\subsection{Component Breakdown}

\begin{center}
\rowcolors{2}{rowgray}{white}
\begin{tabularx}{\textwidth}{L{3.5cm}L{3cm}C{1.5cm}C{2cm}}
\toprule
\rowcolor{primary}
\tableheader{Component} & \tableheader{Dependencies} & \tableheader{Model} & \tableheader{Est. Tokens} \\
\midrule
Wave 0: Foundation schemas & None & Haiku & 15,000 \\
M01: TG/COUM/IR survey & W0 schemas & Sonnet & 40,000 \\
M02: NWW + List analysis & W0 schemas & Opus & 35,000 \\
M03: Current 93 arc & W0, M01 partial & Sonnet & 40,000 \\
M04: Death in June + neofolk & W0, M03 partial & Sonnet + Opus review & 40,000 \\
M05: Coil nexus & All modules + companion study & Opus & 25,000 \\
M06: NIN outlier & M01, M05 & Sonnet + Opus review & 30,000 \\
Network topology synthesis & All modules & Opus & 30,000 \\
Sensitivity review (M04) & M04 draft & Opus (CRAFT-ETHI) & 15,000 \\
Document assembly & All components & Sonnet & 20,000 \\
Verification pass & Assembled doc & Sonnet & 15,000 \\
\midrule
\multicolumn{3}{r}{\textbf{Total estimated:}} & \textbf{305,000} \\
\bottomrule
\end{tabularx}
\end{center}

\subsection{Model Assignment Rationale}

\begin{center}
\rowcolors{2}{rowgray}{white}
\begin{tabularx}{\textwidth}{L{2cm}C{1.5cm}X}
\toprule
\rowcolor{primary}
\tableheader{Model} & \tableheader{Allocation} & \tableheader{Rationale} \\
\midrule
Opus & 30\% & Network topology judgment; ideological analysis (Module 04); NIN outlier framing; synthesis pass; any decisions involving contested claims \\
Sonnet & 60\% & Module surveys; discography tables; document assembly; verification; bibliography formatting \\
Haiku & 10\% & Shared schemas; table templates; scaffolding; source index structure \\
\bottomrule
\end{tabularx}
\end{center}

\section{Wave Execution Plan}

\subsection{Wave Summary}

\begin{center}
\rowcolors{2}{rowgray}{white}
\begin{tabularx}{\textwidth}{L{1.5cm}L{3cm}L{2.5cm}C{2cm}X}
\toprule
\rowcolor{primary}
\tableheader{Wave} & \tableheader{Name} & \tableheader{Models} & \tableheader{Sessions} & \tableheader{Output} \\
\midrule
0 & Foundation & Haiku & 1 & Schemas, templates, source index \\
1 & Parallel Surveys & Sonnet + Opus & 2--3 & Modules 01--04 drafted \\
2 & Synthesis & Opus & 1--2 & Network topology; M05 + M06 complete; sensitivity review \\
3 & Publication & Sonnet & 1 & Final assembled document \\
\bottomrule
\end{tabularx}
\end{center}

\subsection{Wave 0: Foundation}

\begin{center}
\rowcolors{2}{rowgray}{white}
\begin{tabularx}{\textwidth}{L{1.5cm}L{3.5cm}L{2.5cm}X}
\toprule
\rowcolor{primary}
\tableheader{Task} & \tableheader{Description} & \tableheader{Agent} & \tableheader{Output} \\
\midrule
W0-01 & Artist schema template & Haiku & Shared schema: name, era, key albums, key collaborators, network edges \\
W0-02 & Source index setup & Haiku & Indexed bibliography with short-codes for in-text citation \\
W0-03 & Network edge schema & Haiku & Edge definition: artist A $\to$ artist B, type (collab/influence/derivation), evidence \\
W0-04 & Module skeleton files & Haiku & Empty module files with section headers ready for Wave 1 population \\
\bottomrule
\end{tabularx}
\end{center}

\subsection{Wave 1: Parallel Surveys}

\begin{center}
\rowcolors{2}{rowgray}{white}
\begin{tabularx}{\textwidth}{L{1.5cm}L{1.5cm}L{3cm}L{2cm}X}
\toprule
\rowcolor{primary}
\tableheader{Track} & \tableheader{Task} & \tableheader{Description} & \tableheader{Agent} & \tableheader{Output} \\
\midrule
A & W1A-01 & TG/COUM history & Sonnet & Module 01 complete draft \\
A & W1A-02 & Industrial Records catalog & Sonnet (CRAFT-HIST) & Release table with contextual notes \\
B & W1B-01 & NWW + List analysis & Opus & Module 02 complete draft (three-level analysis) \\
B & W1B-02 & NWW List entry survey & Sonnet & Genre/geography distribution table \\
C & W1C-01 & Current 93 arc & Sonnet & Module 03 draft (album-by-album) \\
C & W1C-02 & Death in June + neofolk & Sonnet & Module 04 draft (ideology section flagged for Wave 2 review) \\
\bottomrule
\end{tabularx}
\end{center}

\subsection{Wave 2: Synthesis}

\begin{center}
\rowcolors{2}{rowgray}{white}
\begin{tabularx}{\textwidth}{L{1.5cm}L{3.5cm}L{2.5cm}X}
\toprule
\rowcolor{primary}
\tableheader{Task} & \tableheader{Description} & \tableheader{Agent} & \tableheader{Output} \\
\midrule
W2-01 & Network topology construction & Opus (CRAFT-TOPO) & Complete edge map; 20+ documented relationships \\
W2-02 & Module 04 ideology review & Opus (CRAFT-ETHI) & Reviewed and revised DiJ/neofolk ideology section \\
W2-03 & Module 05: Coil nexus & Opus & Network synthesis; companion study cross-references \\
W2-04 & Module 06: NIN outlier & Opus + Sonnet & Outlier analysis; escape velocity concept developed \\
W2-05 & Cross-module consistency & Sonnet (INTEG) & Entity names, dates, credits consistent across all modules \\
\bottomrule
\end{tabularx}
\end{center}

\subsection{Wave 3: Publication and Verification}

\begin{center}
\rowcolors{2}{rowgray}{white}
\begin{tabularx}{\textwidth}{L{1.5cm}L{3.5cm}L{2.5cm}X}
\toprule
\rowcolor{primary}
\tableheader{Task} & \tableheader{Description} & \tableheader{Agent} & \tableheader{Output} \\
\midrule
W3-01 & Document assembly & Sonnet & Full document with consistent formatting \\
W3-02 & Source verification & Sonnet (VERIFY) & All citations traceable; zero entertainment media \\
W3-03 & Bibliography finalization & Sonnet & Complete, formatted bibliography \\
W3-04 & HITL review gate & CAPCOM & Human approval before final output \\
W3-05 & Final output package & Sonnet & PDF + .tex source + index.html \\
\bottomrule
\end{tabularx}
\end{center}

\subsection{Cache Optimization Strategy}

\begin{itemize}[noitemsep]
  \item Wave 0 outputs (schemas, source index) cached at session start of Wave 1 tracks A, B, C
  \item Artist profiles (Module 01--04) cached for Wave 2 synthesis — 5-minute cache window exploited by starting synthesis immediately after final Wave 1 track completes
  \item Companion Coil discography study loaded once at Wave 2 start and kept in context for Module 05 construction
  \item Network edge map (W2-01 output) cached for Module 05 and 06 construction
\end{itemize}

\subsection{Token Budget}

\begin{center}
\rowcolors{2}{rowgray}{white}
\begin{tabularx}{\textwidth}{L{4cm}C{2.5cm}C{2.5cm}C{2.5cm}}
\toprule
\rowcolor{primary}
\tableheader{Wave} & \tableheader{Opus} & \tableheader{Sonnet} & \tableheader{Haiku} \\
\midrule
Wave 0: Foundation & --- & --- & 15,000 \\
Wave 1A: TG/IR & --- & 80,000 & --- \\
Wave 1B: NWW + List & 35,000 & 10,000 & --- \\
Wave 1C: C93 + DiJ & 10,000 & 80,000 & --- \\
Wave 2: Synthesis & 65,000 & 20,000 & --- \\
Wave 3: Publication & --- & 35,000 & --- \\
\midrule
\textbf{Subtotals} & \textbf{110,000} & \textbf{225,000} & \textbf{15,000} \\
\textbf{Grand Total} & \multicolumn{3}{c}{\textbf{350,000 tokens}} \\
\bottomrule
\end{tabularx}
\end{center}

\subsection{Dependency Graph}

\begin{asciibox}
DEPENDENCY GRAPH

  [W0: Foundation] ---------> [W1A: TG Survey]
       |                           |
       |            -------------> [W1B: NWW + List]
       |            |
       +----------> [W1C: C93 + DiJ]
                            |
                            v
              [W2-01: Network Topology] <-- W1A, W1B, W1C
              [W2-02: Ideology Review] <-- W1C (M04 draft)
              [W2-03: Coil Nexus]     <-- W2-01 + companion study
              [W2-04: NIN Outlier]    <-- W2-01 + W2-03
              [W2-05: Consistency]    <-- W2-01 through W2-04
                            |
                            v
              [W3-01: Document Assembly]
              [W3-02: Source Verification]
              [W3-03: Bibliography]
                            |
              [W3-04: HITL Gate] ---- CAPCOM ---- Human Approval
                            |
              [W3-05: Final Output Package]
\end{asciibox}

\section{Test Plan}

\subsection{Test Categories}

\begin{center}
\rowcolors{2}{rowgray}{white}
\begin{tabularx}{\textwidth}{L{3cm}C{1.5cm}L{2cm}L{2cm}X}
\toprule
\rowcolor{primary}
\tableheader{Category} & \tableheader{Count} & \tableheader{Priority} & \tableheader{Failure} & \tableheader{Purpose} \\
\midrule
Safety-critical & 5 & Mandatory & BLOCK & Non-negotiable source and ideology boundaries \\
Core functionality & 20 & Required & BLOCK & Deliverable acceptance per success criteria \\
Integration & 6 & Required & BLOCK & Cross-module consistency and companion linkage \\
Edge cases & 5 & Best-effort & LOG & Robustness and completeness \\
\bottomrule
\end{tabularx}
\end{center}

\subsection{Safety-Critical Tests}

\begin{center}
\rowcolors{2}{rowgray}{white}
\begin{tabularx}{\textwidth}{L{1.5cm}L{4cm}X}
\toprule
\rowcolor{primary}
\tableheader{Test ID} & \tableheader{Verifies} & \tableheader{Expected Behavior} \\
\midrule
SC-SRC & Source quality & All citations traceable to books, peer-reviewed music journalism, or official release documentation. Zero entertainment blogs, Reddit, or unsourced wikis as primary claims. \\
SC-NUM & Numerical attribution & All sales figures, entry counts, pressing numbers, and dates attributed to specific named sources. \\
SC-ADV & No policy advocacy & Document presents the ideological content of neofolk (Death in June, Boyd Rice) as cultural analysis, not editorial advocacy for or against the artists. \\
SC-ETHI & Ideology attribution accuracy & All claims connecting specific artists to specific political positions cite documented evidence, not scene-wide generalizations. \\
SC-PRON & P-Orridge pronoun accuracy & Genesis P-Orridge referred to with she/her pronouns consistent with P-Orridge's own documented preference in later life. \\
\bottomrule
\end{tabularx}
\end{center}

\subsection{Selected Core and Integration Tests}

\begin{center}
\rowcolors{2}{rowgray}{white}
\begin{tabularx}{\textwidth}{L{1.5cm}L{5cm}X}
\toprule
\rowcolor{primary}
\tableheader{Test ID} & \tableheader{Verifies} & \tableheader{Expected} \\
\midrule
CF-01 & TG arc completeness & Module 01 covers COUM origin through 1981 dissolution with all key albums documented \\
CF-02 & Industrial Records catalog & At least 6 label releases documented beyond TG's own output \\
CF-03 & NWW List three-level analysis & Module 02 addresses artifact, index, and structural map functions \\
CF-04 & Entry count documentation & 291 entries confirmed with expansion from 236 (1979) to 291 (1980) noted \\
CF-05 & C93 album count & At least 8 Current 93 albums analyzed with arc placement \\
CF-06 & Tibet theological arc & Arc from Crowleyan occultism to Christian faith traced with key transition points \\
CF-07 & DiJ folk turn documented & Acoustic guitar introduction traced to specific album (\textit{Burial}, 1984) \\
CF-08 & Ideology section completeness & Module 04 addresses Death's Head imagery, Boyd Rice collaboration, Tibet's distancing \\
CF-09 & Coil nexus edges & Module 05 documents Coil's connections to TG, C93, NWW, DiJ with evidence \\
CF-10 & NIN outlier analysis & Reznor's industrial borrowings and rejection of the label both documented \\
CF-11 & Coil/NIN contact documented & \textit{Fixed} (1992) and \textit{Closer to God} (1994) both analyzed \\
CF-12 & Network edge count & Topology diagram documents at least 20 directed edges with evidence \\
IN-01 & TG $\to$ C93 cascade & Psychic TV connection to Current 93 origin traced across Module 01 and 03 \\
IN-02 & NWW $\leftrightarrow$ C93 consistency & Stapleton's collaborative role described consistently between Module 02 and 03 \\
IN-03 & Companion study linkage & Module 05 contains explicit cross-references to Coil discography study module structure \\
IN-04 & Equinox Event cross-ref & Event appears in both Module 03 (C93 origin) and Module 05 (Coil nexus) consistently \\
IN-05 & World Serpent cross-ref & DiJ distribution infrastructure in Module 04 connects to neofolk network map in Module 06 \\
IN-06 & Self-containment & Reader with no prior knowledge can follow document; all terms defined on first use \\
EC-01 & Disputed facts flagged & NWW List invented-entry controversy handled with both Stapleton and Fothergill positions \\
EC-02 & Temporal currency & Primary reference texts (Keenan 2016 rev.) cited; later scholarship noted where available \\
\bottomrule
\end{tabularx}
\end{center}

\subsection{Verification Matrix}

\begin{center}
\rowcolors{2}{rowgray}{white}
\begin{tabularx}{\textwidth}{XL{3.5cm}C{1.5cm}}
\toprule
\rowcolor{primary}
\tableheader{Success Criterion} & \tableheader{Test ID(s)} & \tableheader{Status} \\
\midrule
1. Origin documentation complete & CF-01, CF-02 & Pending \\
2. NWW List analyzed at all three levels & CF-03, CF-04 & Pending \\
3. Current 93 arc (8+ albums) & CF-05, CF-06 & Pending \\
4. Death in June arc + ideology & CF-07, CF-08 & Pending \\
5. Network topology (20+ edges) & CF-12, IN-01 & Pending \\
6. Coil nexus established & CF-09, IN-03 & Pending \\
7. NIN outlier framing & CF-10 & Pending \\
8. Coil/NIN contact documented & CF-11 & Pending \\
9. Ethical analysis of neofolk & SC-ETHI, CF-08 & Pending \\
10. Sensitivity compliance & SC-SRC, SC-ADV & Pending \\
11. Self-containment & IN-06 & Pending \\
12. Companion study integration & IN-03, IN-04 & Pending \\
\bottomrule
\end{tabularx}
\end{center}

\section{Mission Crew Manifest}

\begin{center}
\rowcolors{2}{rowgray}{white}
\begin{tabularx}{\textwidth}{L{2cm}L{3cm}X}
\toprule
\rowcolor{primary}
\tableheader{Role} & \tableheader{Agent} & \tableheader{Responsibility} \\
\midrule
FLIGHT & Orchestrator & Mission direction; wave go/no-go gates; final integration approval \\
PLAN & Planner & Wave decomposition; parallel track optimization; HITL question generation \\
SCOUT & Research Agent & Supplementary source gathering; gap-fill web research; discography verification \\
EXEC-A & Executor (TG/NWW) & Wave 1A and 1B production; Module 01 and 02 drafts \\
EXEC-B & Executor (C93/DiJ) & Wave 1C production; Module 03 and 04 drafts \\
CRAFT-HIST & History Specialist & Cultural context; era analysis; Industrial Records institutional history \\
CRAFT-TOPO & Topology Specialist & Network edge mapping; relationship graph construction; Module 05 architecture \\
CRAFT-ETHI & Ethics Specialist & Module 04 ideology review; sensitivity compliance; P-Orridge pronoun audit \\
CRAFT-ARCH & Architecture Specialist & Cross-module structure; companion study integration; Module 06 outlier framing \\
VERIFY & Verification Agent & Source validation; cross-module consistency; bibliography completeness \\
INTEG & Integration Agent & IN-series tests; cross-reference validation between all six modules \\
CAPCOM & HITL Interface & Human approval gate at Wave 2/3 boundary; only channel crossing to user \\
BUDGET & Resource Manager & Token budget tracking per wave; alert at 80\% consumption \\
SURGEON & Health Monitor & Research quality degradation detection; escalation to FLIGHT \\
TOPO & Topology Manager & Agent team restructuring if mid-mission scope change required \\
RETRO & Retrospective & Post-wave analysis; lessons captured for companion and future missions \\
LOG & Chronicle Agent & Commit messages; milestone summaries; audit trail \\
\bottomrule
\end{tabularx}
\end{center}

\section{Execution Summary}

\begin{center}
\rowcolors{2}{rowgray}{white}
\begin{tabularx}{\textwidth}{L{4cm}X}
\toprule
\rowcolor{primary}
\tableheader{Metric} & \tableheader{Value} \\
\midrule
Activation profile & Fleet (17 roles) \\
Total modules & 6 \\
Parallel tracks (Wave 1) & 3 \\
Estimated total tokens & 305,000--350,000 \\
Estimated sessions & 4--5 \\
Context windows & 14--18 \\
Primary model & Sonnet (60\%) \\
Judgment model & Opus (30\%) \\
Scaffold model & Haiku (10\%) \\
Safety-critical test count & 5 \\
Total test count & 36 \\
HITL gates & 1 (Wave 2/3 boundary) \\
Companion study dependency & Coil discography study (Module 05 cross-reference) \\
\bottomrule
\end{tabularx}
\end{center}

\vspace{2em}

\begin{quotebox}
\textit{``Categories strain, crack and sometimes break, under their burden ---\\
step out of the space provided.''}

\smallskip
--- Nurse With Wound List, opening declaration (1979)

\smallskip
The network this study maps was built by people who understood that the spaces between artists
are where the actual intelligence lives. Not the nodes. The edges. The collaborations, the
departures, the deliberate refusals. The Amiga Principle, applied at cultural scale: small,
principled components, faithfully iterated, producing staggering complexity from the
connections between them. That is what the post-industrial network was. That is why it matters.
\end{quotebox}

\end{document}
